%%%%%%%%%%%%%%%%
%%% deut 17 %%%
%%%%%%%%%%%%%%%%

\Chapter{17}

\Verse{1}
{
	\hDtn{לֹא תִזְבַּח לַיהֹוָה אֱלֹהֶיךָ שׁוֹר וָשֶׂה אֲשֶׁר יִהְיֶה בוֹ מוּם כֹּל דָּבָר רָע כִּי תוֹעֲבַת יְהֹוָה אֱלֹהֶיךָ הוּא׃}
}
{
	\eDtn{
		“You shall not sacrifice to YHWH, your {\God}, a bull or a
		sheep in which will be any injury, any bad thing, because
		that is an offensive thing of YHWH, your {\God}.
	}
}
{
}

\Verse{2}
{
	\hDtn{כִּי יִמָּצֵא בְקִרְבְּךָ בְּאַחַד שְׁעָרֶיךָ אֲשֶׁר יְהֹוָה אֱלֹהֶיךָ נֹתֵן לָךְ אִישׁ אוֹ אִשָּׁה אֲשֶׁר יַעֲשֶׂה אֶת הָרַע בְּעֵינֵי יְהֹוָה אֱלֹהֶיךָ לַעֲבֹר בְּרִיתוֹ׃}
}
{
	\eDtn{
		“If there will be found among you, in one of your gates
		that YHWH, your {\God}, is giving you, a man or woman who
		will do what is bad in the eyes of YHWH, your {\God}, to
		violate His covenant,
	}
}
{
}

\Verse{3}
{
	\hDtn{וַיֵּלֶךְ וַיַּעֲבֹד אֱלֹהִים אֲחֵרִים וַיִּשְׁתַּחוּ לָהֶם וְלַשֶּׁמֶשׁ אוֹ לַיָּרֵחַ אוֹ לְכל צְבָא הַשָּׁמַיִם אֲשֶׁר לֹא צִוִּיתִי׃}
}
{
	\eDtn{
		and will go and serve other gods and bow to them and to
		the sun or to the moon or to any of the array of the
		skies, which I did not command,
	}
}
{
}

\Verse{4}
{
	\hDtn{וְהֻגַּד לְךָ וְשָׁמָעְתָּ וְדָרַשְׁתָּ הֵיטֵב וְהִנֵּה אֱמֶת נָכוֹן הַדָּבָר נֶעֶשְׂתָה הַתּוֹעֵבָה הַזֹּאת בְּיִשְׂרָאֵל׃}
}
{
	\eDtn{
		and it will be told to you and you will hear it, and you
		will inquire well, and, here, it is true, the thing is
		right, this offensive thing has been done in Israel,
	}
}
{
}

\Verse{5}
{
	\hDtn{וְהוֹצֵאתָ אֶת הָאִישׁ הַהוּא אוֹ אֶת הָאִשָּׁה הַהִוא אֲשֶׁר עָשׂוּ אֶת הַדָּבָר הָרָע הַזֶּה אֶל שְׁעָרֶיךָ אֶת הָאִישׁ אוֹ אֶת הָאִשָּׁה וּסְקַלְתָּם בָּאֲבָנִים וָמֵתוּ׃}
}
{
	\eDtn{
		then you shall bring that man or that woman who did this
		bad thing out to your gates, the man or the woman, and you
		shall stone them with stones so they die.
	}
}
{
}

\Verse{6}
{
	\hDtn{עַל פִּי שְׁנַיִם עֵדִים אוֹ שְׁלֹשָׁה עֵדִים יוּמַת הַמֵּת לֹא יוּמַת עַל פִּי עֵד אֶחָד׃}
}
{
	\eDtn{
		“On the word of two witnesses or three witnesses shall the
		one who is to die be put to death. One shall not be put to
		death on the word of one witness.
	}
}
{
}

\Verse{7}
{
	\hDtn{יַד הָעֵדִים תִּהְיֶה בּוֹ בָרִאשֹׁנָה לַהֲמִיתוֹ וְיַד כּל הָעָם בָּאַחֲרֹנָה וּבִעַרְתָּ הָרָע מִקִּרְבֶּךָ׃}
}
{
	\eDtn{
		The witnesses’ hand shall be on him first to put him to
		death, and all the people’s hand after that. So you shall
		burn away what is bad from among you.
	}
}
{
%	\footnote{
%		witnesses first … all the people after. Numerous commentators discuss
%		the first half of this verse but not the second. It is of course
%		significant that witnesses must know that, if they testify against
%		someone in a capital case, they will have to strike the first blows
%		against the defendant themselves. That conveys vividly how serious their
%		testimony is. But it is at least as significant that all the people
%		participate in the execution as well. Capital punishment is the act of a
%		community—in ancient Israel and even more so in democracies. In biblical
%		law, the community must know that if they are going to have capital
%		punishment they must understand that this means that the community as a
%		whole is taking lives. So that there can be no mistaking this point, the
%		entire community must participate in the execution. In postbiblical
%		societies where the entire people does not participate personally, this
%		point may easily be forgotten, but it is true all the same: a community
%		that has capital punishment is responsible as a community for the lives
%		it takes. It is not just the executioner or just the witnesses or just
%		the courts. Perhaps this fact—that in ancient Israel all the people had
%		to carry out the execution—was one of the forces that led Israel to make
%		capital punishment almost impossible in rabbinic law.
%	}
}

\Verse{8}
{
	\hDtn{כִּי יִפָּלֵא מִמְּךָ דָבָר לַמִּשְׁפָּט בֵּין דָּם לְדָם בֵּין דִּין לְדִין וּבֵין נֶגַע לָנֶגַע דִּבְרֵי רִיבֹת בִּשְׁעָרֶיךָ וְקַמְתָּ וְעָלִיתָ אֶל הַמָּקוֹם אֲשֶׁר יִבְחַר יְהֹוָה אֱלֹהֶיךָ בּוֹ׃}
}
{
	\eDtn{
		“If a matter for judgment will be too daunting for you,
		between blood and blood, between law and law, and between
		injury and injury, matters of disputes in your gates, then
		you shall get up and go up to the place that YHWH, your
		{\God}, will choose.
	}
}
{
%	\footnote{
%		If a matter for judgment will be too daunting for you. In every law code
%		there must be a mechanism for change and for application to new and
%		difficult situations. Even if the law is divine law, there must be such
%		a mechanism, and the Torah recognizes this and provides for it. It
%		directs that in such difficult questions, the authorities (judges and
%		priests; i.e., authorities in law and religion) in each age shall
%		determine what to do. This has always been done in Judaism. The most
%		obvious distinction among the movements in Judaism has been their
%		different views of how the law changes. The primary consideration when
%		authorities, including scholars and rabbis, determine that a law is
%		changed is that they do so with wisdom and reverence, and not with
%		arrogance. Footnote 17:8. blood and blood. Cases involving the taking of
%		life. Footnote 17:8. between law and law. Cases that are complex because
%		more than one law applies. Footnote 17:8. between injury and injury.
%		Cases involving physical hurt. Footnote 17:8. disputes. The term (Hebrew
%		rîb) meant a quarrel, as in the quarreling at the rock that leads to the
%		name Meribah (Exodus 17; Numbers 20), and it came to refer to a lawsuit,
%		a legal dispute. Israel’s prophets use it to describe God’s judgment of
%		Israel in legal terms (Hos 12:3; Mic 6:2). This is known in biblical
%		scholarship as a “covenant lawsuit.” The covenant between God and Israel
%		is a legal relationship, formulated in the legal terminology of the
%		ancient Near East (see the comment on Exod 20:1). In a rîb oracle, the
%		prophet describes God as bringing a case against His people for
%		violating the covenant, which can end in judgment and punishment.
%	}
}

\Verse{9}
{
	\hDtn{וּבָאתָ אֶל הַכֹּהֲנִים הַלְוִיִּם וְאֶל הַשֹּׁפֵט אֲשֶׁר יִהְיֶה בַּיָּמִים הָהֵם וְדָרַשְׁתָּ וְהִגִּידוּ לְךָ אֵת דְּבַר הַמִּשְׁפָּט׃}
}
{
	\eDtn{
		And you shall come to the Levite priests and to the judge
		who will be in those days, and you shall inquire, and they
		will tell you the matter of judgment.
	}
}
{
%	\footnote{
%		the Levite priests. Historically, originally all Levites were priests.
%		Later the Levite group that identified itself as descendants of Aaron
%		took exclusive hold of the priesthood and limited all the other Levites
%		to a secondary role. Thereafter a distinction was made between priests
%		and Levites. Deuteronomy reflects the original status of all Levites and
%		does not make this distinction between priests and Levites. (I have
%		discussed this history of priesthood in Who Wrote the Bible?)
%	}
}

\Verse{10}
{
	\hDtn{וְעָשִׂיתָ עַל פִּי הַדָּבָר אֲשֶׁר יַגִּידוּ לְךָ מִן הַמָּקוֹם הַהוּא אֲשֶׁר יִבְחַר יְהֹוָה וְשָׁמַרְתָּ לַעֲשׂוֹת כְּכֹל אֲשֶׁר יוֹרוּךָ׃}
}
{
	\eDtn{
		And you shall do according to the word on the matter that
		they will tell you from that place that YHWH will choose,
		and you shall be watchful to do according to everything
		that they will instruct you.
	}
}
{
}

\Verse{11}
{
	\hDtn{עַל פִּי הַתּוֹרָה אֲשֶׁר יוֹרוּךָ וְעַל הַמִּשְׁפָּט אֲשֶׁר יֹאמְרוּ לְךָ תַּעֲשֶׂה לֹא תָסוּר מִן הַדָּבָר אֲשֶׁר יַגִּידוּ לְךָ יָמִין וּשְׂמֹאל׃}
}
{
	\eDtn{
		You shall do it according to the word of the instruction
		that they will give you and according to the judgment that
		they will say to you. You shall not turn from the thing
		that they will tell you, right or left.
	}
}
{
}

\Verse{12}
{
	\hDtn{וְהָאִישׁ אֲשֶׁר יַעֲשֶׂה בְזָדוֹן לְבִלְתִּי שְׁמֹעַ אֶל הַכֹּהֵן הָעֹמֵד לְשָׁרֶת שָׁם אֶת יְהֹוָה אֱלֹהֶיךָ אוֹ אֶל הַשֹּׁפֵט וּמֵת הָאִישׁ הַהוּא וּבִעַרְתָּ הָרָע מִיִּשְׂרָאֵל׃}
}
{
	\eDtn{
		And the man who will act presumptuously, not listening to
		the priest who is standing to serve YHWH, your {\God}, there,
		or to the judge: that man shall die. So you shall burn
		away what is bad from Israel.
	}
}
{
}

\Verse{13}
{
	\hDtn{וְכל הָעָם יִשְׁמְעוּ וְיִרָאוּ וְלֹא יְזִידוּן עוֹד׃}
}
{
	\eDtn{
		And all the people will listen and fear and won’t act
		presumptuously anymore.
	}
}
{
}

\Verse{14}
{
	\hDtn{כִּי תָבֹא אֶל הָאָרֶץ אֲשֶׁר יְהֹוָה אֱלֹהֶיךָ נֹתֵן לָךְ וִירִשְׁתָּהּ וְיָשַׁבְתָּה בָּהּ וְאָמַרְתָּ אָשִׂימָה עָלַי מֶלֶךְ כְּכל הַגּוֹיִם אֲשֶׁר סְבִיבֹתָי׃}
}
{
	\eDtn{
		“When you’ll come to the land that YHWH, your {\God}, is
		giving you, and you’ll take possession of it and live in
		it, and you’ll say, ‘Let me set a king over me like all
		the nations that are around me,’
	}
}
{
}

\Verse{15}
{
	\hDtn{שׂוֹם תָּשִׂים עָלֶיךָ מֶלֶךְ אֲשֶׁר יִבְחַר יְהֹוָה אֱלֹהֶיךָ בּוֹ מִקֶּרֶב אַחֶיךָ תָּשִׂים עָלֶיךָ מֶלֶךְ לֹא תוּכַל לָתֵת עָלֶיךָ אִישׁ נכְרִי אֲשֶׁר לֹא אָחִיךָ הוּא׃}
}
{
	\eDtn{
		you shall set a king over you whom YHWH, your {\God}, will
		choose! You shall set a king from among your brothers over
		you. You may not put a foreign man, who is not your
		brother, over you.
	}
}
{
%	\footnote{
%		you shall set a king over you. The Torah permits Israel to have a king
%		if they wish. It is placed in their hands to decide if they want one:
%		“When … you’ll say, ‘Let me set a king over me … ,’ you shall set a king
%		over you.” But the king must be chosen by God (see the next comment).
%		Later, Gideon and Samuel will question whether it is right for Israel to
%		have a king, saying that God should be Israel’s king (Judg 8:23; 1 Sam
%		8:6). But they do not base their opposition to monarchy on the Torah.
%		Both Gideon and Samuel are judges and also priests (Judg 8:24–27), and
%		they naturally think of God’s will as being properly expressed through
%		priests, and so they see kings as a threat to the order that they know.
%		But the Torah permits the people to have kings if they wish. And so God
%		tells Samuel to give the people a king, even though their wish for a
%		king, in that particular case, involves a rejection of God (1 Sam
%		8:7,22). Footnote 17:15. whom God will choose. This means that a prophet
%		must designate the king as chosen by God. Thus Samuel designates Saul
%		and David, Nathan designates Solomon, and Ahijah designates Jeroboam.
%		The stories of the kings of Israel and Judah in the books of Samuel,
%		Kings, and Chronicles convey that kings are acceptable if the people
%		want them, but their power is limited and balanced by the power of
%		prophets. And kings who do not listen to prophets, or who imprison or
%		execute them, do so at their peril.
%	}
}

\Verse{16}
{
	\hDtn{רַק לֹא יַרְבֶּה לּוֹ סוּסִים וְלֹא יָשִׁיב אֶת הָעָם מִצְרַיְמָה לְמַעַן הַרְבּוֹת סוּס וַיהֹוָה אָמַר לָכֶם לֹא תֹסִפוּן לָשׁוּב בַּדֶּרֶךְ הַזֶּה עוֹד׃}
}
{
	\eDtn{
		Only, he shall not get himself many horses, and he shall
		not bring the people back to Egypt in order to get many
		horses, when YHWH has said to you, ‘You shall not go back
		this way ever again.’
	}
}
{
%	\footnote{
%		not get himself many horses. Maintaining forces with a great many horses
%		means the development of a small upper class plus a large class of serfs
%		or servants to support and serve them. This is one of several laws in
%		the Torah that work to counter the development of such a two-tiered
%		feudal system. And so, when King David defeats Hadadezer and captures
%		1,700 chariots and horsemen, he hamstrings all but a small number of the
%		horses (2 Sam 8:4). Why would he deliberately give up a prize of such
%		tremendous value? Because he is obeying this Law of the King. Footnote
%		17:16. YHWH has said to you, ‘You shall not go back.’ No such words from
%		God are reported prior to this. This may appear to be a problem, but it
%		may simply be that this is Moses’ report of God’s words on this matter.
%	}
}

\Verse{17}
{
	\hDtn{וְלֹא יַרְבֶּה לּוֹ נָשִׁים וְלֹא יָסוּר לְבָבוֹ וְכֶסֶף וְזָהָב לֹא יַרְבֶּה לּוֹ מְאֹד׃}
}
{
	\eDtn{
		And he shall not get himself many wives, so his heart will
		not turn away. And he shall not get himself very much
		silver and gold.
	}
}
{
}

\Verse{18}
{
	\hDtn{וְהָיָה כְשִׁבְתּוֹ עַל כִּסֵּא מַמְלַכְתּוֹ וְכָתַב לוֹ אֶת מִשְׁנֵה הַתּוֹרָה הַזֹּאת עַל סֵפֶר מִלִּפְנֵי הַכֹּהֲנִים הַלְוִיִּם׃}
}
{
	\eDtn{
		And it will be, when he sits on his kingdom’s throne, that
		he shall write himself a copy of this instruction on a
		scroll from in front of the Levite priests.
	}
}
{
%	\footnote{
%		a copy of this instruction. It is unclear whether this means a copy of
%		this Law of the King or a copy of the full law code of Deuteronomy in
%		which it is now contained. In either case, this means that Israel has a
%		constitutional monarchy. Baruch Halpern has shown substantial evidence
%		that Israel’s kings were bound by this Law of the King from the time of
%		the first king, Saul (Halpern, The Constitution of the Monarchy in
%		Israel). That means that Israel was historically a constitutional
%		monarchy, with a written constitution, more than a thousand years B.C.E.
%		Footnote 17:18. a copy … from in front of the Levite priests. That is,
%		the Levites have the text of the law, and the king makes himself a copy
%		from the text that they have before them.
%	}
}

\Verse{19}
{
	\hDtn{וְהָיְתָה עִמּוֹ וְקָרָא בוֹ כּל יְמֵי חַיָּיו לְמַעַן יִלְמַד לְיִרְאָה אֶת יְהֹוָה אֱלֹהָיו לִשְׁמֹר אֶת כּל דִּבְרֵי הַתּוֹרָה הַזֹּאת וְאֶת הַחֻקִּים הָאֵלֶּה לַעֲשֹׂתָם׃}
}
{
	\eDtn{
		And it shall be with him, and he shall read it all the
		days of his life, so that he will learn to fear YHWH, his
		{\God}, to observe all the words of this instruction and
		these laws, to do them,
	}
}
{
}

\Verse{20}
{
	\hDtn{לְבִלְתִּי רוּם לְבָבוֹ מֵאֶחָיו וּלְבִלְתִּי סוּר מִן הַמִּצְוָה יָמִין וּשְׂמֹאול לְמַעַן יַאֲרִיךְ יָמִים עַל מַמְלַכְתּוֹ הוּא וּבָנָיו בְּקֶרֶב יִשְׂרָאֵל׃}
}
{
	\eDtn{
		so his heart will not be elevated above his brothers, and
		so he will not turn from the commandment, right or left,
		and so he will extend days over his kingdom, he and his
		sons, within Israel.
	}
}
{
}
